\section{Two partial results}\label{sec:partial}

\subsection{Effective free-metabelian growth (G9)}\label{sec:g9}
\entry{G9}{Collection entry, with no individual author listed}
The question asks for the growth rate of a free metabelian group of
rank $r$, with its standard generators. An approximation algorithm and
an interval leave the exact constant undetermined.

Let $V_r(n)$ be the cardinality of the radius-$n$ ball in $M_r$ and
$\lambda_r=\lim_n V_r(n)^{1/n}$. Put
\[
 P(t)=\#\{s_1>\cdots>s_k>0:\ \sum s_i\le t\},\qquad
 K_n=(n+1)^3P(n+1)^2,
\]
including the empty list in $P(t)$.

\begin{theorem}\label{prop:g9}
Uniformly in $r$ and $n\ge1$,
\begin{equation}\label{eq:growth-bounds}
 \max\left\{1,\left(\frac{V_r(n)}{K_n}\right)^{1/(n+2)}\right\}
 \le\lambda_r\le V_r(n)^{1/n}.
\end{equation}
For these endpoints $L_n,U_n$,
\begin{equation}\label{eq:growth-modulus}
 \log(U_n/L_n)\le
 \frac{2\log(2r+1)+3\log(n+1)+4\sqrt{n+1}}{n+2}.
\end{equation}
Thus $\lambda_r$ is uniformly computable. The retained rank-two
certificates give
\[
 \boxed{2.676891785\le\lambda_2\le2.943737759.}
\]
\end{theorem}

\begin{proof}[Proof of the approximation assertion]
A word in $F_r$ is a path in $\Z^r$. Its finitely supported net
integral edge flow determines its image in $M_r$: this is the classical
flow criterion, obtained from the first homology of the abelian cover
\cite{guba,mruv}. In particular $V_r(n)$ is computable by enumerating
words of length at most $n$ and deduplicating flows. No shortest-word
oracle is required.

Use the first coordinate as height. A strict bridge path starts at
height zero, is positive thereafter, and ends at its maximum height $H$.
Its edges, indexed by their upper height, lie in the band $(0,H]$.
Concatenating bridge flows of equal height places them in disjoint
bands. Restricting the product flow to these bands and recovering
successive endpoints from boundaries recovers every factor. If
$b(n,H)$ counts bridge elements with a representative of length at
most $n$ and height $H$, then
\[
 b(n,H)^k\le V_r(kn),\qquad
 b(n):=\sum_{H=1}^n b(n,H)\le n\lambda_r^n.
\]

For each element with a strict half-space representative, choose one
deterministically. Cut it at the last highest remaining vertex, then
the last lowest, and so on. The positive successive spans strictly
decrease and their sum is at most its length. Reflect downward pieces
and concatenate them into a bridge. Given its \emph{flow} and the
span list, restrictions to disjoint bands recover the original piece
flows and their endpoints; reversing reflections and translations
recovers the original element. Thus, writing $h(n)$ for the number
of half-space elements,
\[
 h(n)\le P(n)b(n)\le P(n)n\lambda_r^n.
\]
This adapts the classical bridge unfolding argument~\cite{hw}; the
additional requirement here is recovery from total flow rather than
from a path, allowing repeated edges and cancellations.

Split an arbitrary representative $w=UV$ of length at most $n$ at a
vertex of globally minimal height, after $i$ letters. Then
$a_1U^{-1}$ and $a_1V$ are strict half-space words of lengths at most
$i+1$ and $n-i+1$, and their quotient recovers $w$. Hence
\[
 V_r(n)\le\sum_{i=0}^n h(i+1)h(n-i+1)
 \le(n+1)^3P(n+1)^2\lambda_r^{n+2}.
\]
This proves the lower bound. Submultiplicativity of balls gives the
upper bound. For $s>0$,
\[
 P(t)\le e^{st}\prod_{j\ge1}(1+e^{-sj})
 \le\exp(st+1/s),
\]
so $P(t)\le e^{2\sqrt t}$. Combining this with
$V_r(n)\le(2r+1)^n$ gives~\eqref{eq:growth-modulus}. The effective
error bound and exact finite counts yield rational outward approximations
of any requested precision. The procedure need not be practical.
\end{proof}

The numerical lower bound uses a free monoid of bridge atoms, with
infinite completed boundary and shear families. The upper bound uses
648 verified rewriting relations, a 1,968-state avoidance automaton
and exact positive-vector inequalities. Appendix~\ref{app:g9-certificate}
gives the finite certificate structure and the final rational series.
The flow-support connection construction used to shorten representatives
is prior work~\cite[Section 2.7]{mruv}; no new general geodesic algorithm
or current best-in-literature bound is asserted. The method's recorded
extension to free solvable groups is not another problem count.
\evidence{G9}{flow-growth-proof.md}

\subsection{Infinitely many special five-strand braids (B9)}\label{sec:b9}
\entry{B9}{P.~Dehornoy}
Let $S(\sigma_i)=\sigma_{i+1}$ and
\[
 a\mathbin{\triangleright}b=aS(b)\sigma_1S(a)^{-1}.
\]
The special braids are the closure of $1$ under this operation in
$B_\infty$. How many lie in each standard subgroup $B_N$?
Intermediate terms in a special expression may have more strands
than its final value.

\begin{theorem}\label{prop:b9}
For every $n\ge3$ the braid
\begin{align*}
 \beta_n={}&\sigma_1^n\sigma_3^{n-2}\sigma_2^{1-n}
 \sigma_3^2\sigma_4^{-1}\sigma_2\sigma_1
 \sigma_3^{n-1}\sigma_4^{2-n}\sigma_2^{-n}
\end{align*}
is special, has standard strand support exactly five, and the
$\beta_n$ are pairwise distinct. There are therefore countably
infinitely many special braids in every $B_N$ with $N\ge5$.
\end{theorem}
\begin{proof}
Define $u_0=1$, $u_1=\sigma_1$, and
$u_{n+1}=u_n\triangleright u_{n-1}$. Induction gives
$u_nS(u_{n-1})=\sigma_1^n$. For $n\ge3$ set
\[
 q_n=u_{n-3},\qquad A_n=\sigma_1^n\sigma_3^{n-2}\sigma_2^{1-n}.
\]
Expanding the recurrence three times and commuting distant generators
gives $u_nS^3(q_n)=A_n$. Now take the explicitly special expression
\[
 v_n=1\triangleright((1\triangleright q_n)\triangleright1),
 \qquad \beta_n=u_n\triangleright v_n.
\]
Expansion yields
\[
 v_n=S^2(q_n)\sigma_2^2\sigma_3^{-1}S^3(q_n)^{-1}\sigma_1.
\]
In the next operation, move $S^4(q_n)^{-1}$ past $\sigma_2\sigma_1$
and apply the recurrence. The result is
\[
 \beta_n=A_n\sigma_3^2\sigma_4^{-1}\sigma_2\sigma_1S(A_n)^{-1},
\]
which is the displayed five-strand word, of length $6n-1$ as written
and exponent sum three.

Use the unreduced Burau representation at $t=-1$, whose nontrivial
two-by-two block is
\[
 T=\begin{pmatrix}2&-1\\1&0\end{pmatrix}=I+N,
 \qquad N^2=0,\quad T^k=I+kN\quad(k\in\Z).
\]
Multiplying the fifth row through the displayed blocks gives
\[
 (0,\ n(n-1),\ 1-n^2,\ -n^2+3n-2,\ n^2-2n+2).
\]
Its second entry is strictly increasing for $n\ge3$ and nonzero,
whereas every standard four-strand braid has fifth row $e_5$.
This proves distinctness and strand support; faithfulness is not
needed to distinguish unequal images. The same representation extends
by identities for $N\ge5$. Countability of $B_N$ finishes the proof.
\end{proof}

The small-strand counts $1,2,4$ follow from Dehornoy's prior structural
results~\cite[Sections 4--5]{dehornoy}. In particular a special braid
in $B_N$ has exponent sum between $0$ and $N-1$; sum zero gives $1$
and sum $N-1$ gives $\tau_{N-1}=\sigma_{N-1}\cdots\sigma_1$.
The exponent-one elements are exactly $a\triangleright1$ with
$a$ special in $B_{N-1}$, and this map is injective. Thus
\[
 \mathcal C\cap B_3=\{1,\sigma_1,\sigma_1^2\sigma_2^{-1},
 \sigma_2\sigma_1\}.
\]
In $B_4$ only exponent two remains. Writing
$I_p(A)=A\tau_p S(A)^{-1}$, that sector corresponds to right cosets
$AB_2$ in $B_3$ admitting $A=aS(c)$ with both $a,c$ special.
Four known cosets give at least ten special braids in total; exhaustion
was not proved.

Later reductions decide the complete partner set for each supplied
first braid and exclude several infinite first-input families. For
example, for $I_1(\beta_n)$, $n\ge3$, an invariant Burau entry
$n(n-1)$ excludes every partner returning to the required three-strand
subgroup. These results do not classify unrestricted first inputs.
The global four-strand question therefore remains open here.
\evidence{B9}{infinite-family-proof.md}
